%=============================================================================
% REPRODUCTION
%=============================================================================
\section{Reproducing the Benchmark}
\label{sec:repro}

\cE{We ran the reference implementation~\cite{weng2024graph} unmodified---their
code, their data, their protocol---on a pinned software stack, and recovered their
published numbers to within \textbf{7.3\%} across twenty comparisons, with AAGCN
and DCRNN matching to four decimal places. Three deviations from their environment
were forced and are documented with the exact error that forces each; none changes
a reported figure. The published results are real.}

\cE{What they measure is a different question, and three properties of the
reference evaluation account for the entire apparent advantage.}

\cE{\textbf{The reported column includes training data.} The column labelled
``cross validated'' is computed on a loader spanning the full series rather than
a held-out split; it is a fit statistic, and the column names are effectively
inverted relative to their meanings. \textbf{RMSE is averaged per window, not
pooled.} By Jensen's inequality that is systematically lower, and on this target
the gap is $\sim$39\%. \textbf{The naive baseline is absent.} Adding persistence
to their own reported column reverses the result: it beats \emph{every} model
there, on both MAE and RMSE. An independent re-implementation sharing no code
with either codebase reproduces this, so it is not an artefact of one.}

\cF{Published advantages here therefore do not survive contact with a naive
baseline on held-out data. We take only the \emph{architectures} forward, under
\S\ref{sec:setup}; numbers below are comparable to each other, not to the source
table.}
